\chapter{Method}


\section{Parallel Louvain}
There has been relatively little work on parallelising Louvain, largely because the algorithm is non trivial to parallelise. Louvain is non trivial to parallelise as a result of the constantly changing parameters used in modularity evaluation, the usage of a lock to safeguard these parameters degrade the algorithm into a potentially slower sequential algorithm. Eventually consistent parameters being used in modularity evaluation still caries the same problems as a parallel version relying on stale data, although to a reduced level. An important parallel enabling modification to Louvain is discussed by Lu in \cite{lu2015parallel}, this implementation proves the problem of negative gain and addresses this concern in an effective manner. Importantly, it provides one easy step of parallelisation for us that does not cause negative gain and that is the vertex following heuristic they introduced. The work by Lu proved that not only were they able to obtain significant speedups over the serial algorithm, they were also able to produce communities with higher modularity. 

\subsection{Negative gain scenario}
Although the Louvain modulairty algorithm always improves modularity in its serial implementation, this property is not maintained in parallel implementations. Unfortunately this means that it is possible for the parallel implementation to worsen the modularity of a graph and in this case steps must be made to prevent or mitigate this issue. This negative gain scenario present in Parallel Louvain is discussed and proven by Lu \cite{lu2015parallel}.



\section{Dataflow Louvain}
As far as the author is aware, this is the only literature present that addresses how the Louvain modularity optimisation algorithm could be modified such that it is capable of running as a dataflow program. The potential pitfalls with parallelising Louvain modularity optimisation is present with an even greater effect when considering a dataflow implementation of the algorithm. This is ultimately due to a simple but obvious limitation of most dataflow programming, dataflow programming pulls in its data and it is unable to push and propagate changes immediately unlike a parallel implementation which could in the very least place a bound on the staleness of the data. This is an area of concern when considering any modularity optimisation algorithm since it requires changing data to update the various parameters used in determining modularity such as $k_i$, $\sum_{tot}$ and other variables, this is as a result of these algorithms optimising for a global metric ($Q$), a better solution to parallelise graph clustering is to obtain some local metric to optimise, unfortunately, to the authors knowledge, no such method exists in published literature. 

\subsection{Non-termination}
Unfortunately as a result of our computations being performed on stale data, it is possible to end up in a deadlock. This case is trivial to demonstrate, consider the simple case where there exists a graph of size two. Let's state that when performing the version of Louvain without updates and relying on stale data there is a $\Delta{Q} > 0$ at both vertices, this means that both will perform a move to each others community. Following this move, it becomes clear that the next step could also simply be to repeat this swap, this means that some method is needed to prevent non-termination. 

\subsection{Counterfactual swapping}
Given our usage of differential dataflow, we are able to mitigate some non-termination behaviour relatively cheaply. When we are about to make a decision to move to another community, we do this by examining the modularity gain of moving to the vertex of one of our edges. We can simply perform a join to determine where our neighbour would move, but more importantly we are then also able to evaluate the modularity gain for both of these operations and only perform the community transfer which maximises global modularity. While this may seem counter-intuitive since we limit the community transfers made, it mitigates both non-termination negative gain scenarios to a certain extent, unfortunately it does not however prevent non-termination nor negative-gain. 
The swapping technique above does not prevent non-termination. It is able to mitigate non-termination when considering two vertices however, when the number of vertices that perform cyclic swaps increase, it is obviously unable to account for such a scenario. Consider the case of a set of vertices where each vertex moves in a circular repeating pattern between communities, lets label these vertices as $i,j,k$ such that $i \in C(j),j \in C(j), k \in C(k)$ and $C(i) \neq C(j) \neq C(k)$. Consider the case where $i$ considers moving to $C(j)$, $j$ considers moving to $C(k)$ and $k$ considers moving to $C(i)$ simultaneously. In this case these vertices will see the global modularity as being $Q_{new_i} = Q_{old} + \Delta{Q_{i \rightarrow C(j)}}$, $Q_{new_j} = Q_{old} + \Delta_{Q_{j \rightarrow C(k)}}$ and $Q_{new_k} = Q_{old} + \Delta{Q_{k \rightarrow C(i)}}$ respectively. The critical error here is obvious, these new values of $Q$ are completely wrong and could potentially lead to non-termination, the real value of $Q$ is in fact given by $Q_{new} = Q_{old} + \Delta{Q_{i \rightarrow C(j)}} + \Delta{Q_{j \rightarrow C(k)}} + \Delta{Q_{k \rightarrow C(i)}}$. Note how in the above formula a determination of one transfer directly impacts the values used in the computation of $\Delta{Q}$ in other transfers, this is why Louvain is traditionally a sequential algorithm. 

\subsection{Recognising where deadlocks (non-termination) may occur}
Ultimately deadlocks occur when vertices continuously move around its neighbours and are also able to move back to their initial community structure, following this revelation it becomes easy to recognise that this implies that deadlocks occur in connected components although further work should be attempted to prove if this is exactly the case. A simple yet naive approach to this problem could be to simply agglomerate connected components into a single structure, however this is not always a viable approach as there exists graphs with the number of nodes in communities being smaller than the number of nodes in connected components, as a result using this approach would simply break our community detection algorithm for some graphs. As a result of our graph dataset containing graphs such that the size of one of its strongly connected components was the same size of the graph itself, this was obviously not a practical option. 


\subsection{Implementation}
Ultimately, our method is not dissimilar to the traditional Louvain \cite{blondel2008fast} algorithm for modularity optimisation. The algorithm is expressed as a dataflow program, however the equation to measure both modularity and modularity gain remains the same as the traditional algorithm. One fundamental restriction that had to be made in our implementation however, is to order our nodes, this was quite simply achieved by assigning an integer to every node, this is ultimately due to implementation details of the open source library that was used to implement our algorithm. Like the initial Louvain algorithm, this modification is able to scale up to networks that measure in the millions of edges and is only limited by the memory it consumed. It is important to note that Differential Dataflow implements logic simply by running it until the output data converges to a fixed point. Our algorithm only computes the first phase and the Louvain algorithm is applied repeatedly until a fixed point is reached, modifications such as custom stopping criteria can still be made through simply applying the identity function to the various mappings to force convergence even when the stopping criteria is probabilistic. The algorithm works by first assigning a unique community to every node in the graph, that is, the number of nodes and communities are equal at the start of the first iteration. For each node $i$ in the graph, we consider the neighbours $j$ of $i$ and obtain the modularity gain ($\Delta Q$) of moving it from its community and to the community of its neighbour, when performing this we need to keep track of the highest gain in modularity. After we have determined the highest gain in modularity, we perform the act of moving the node over to its neighbours community if $\Delta Q > 0$, otherwise the node remains in its current community. 
\subsubsection{Justification for not providing agglomeration}
As noted previously, we do not perform any agglomeration in our Louvain based method, the reasoning for this is simple, it does not make sense to agglomerate the clusters when using Louvain in a stream-like context. In addition to this, new nodes being added will not carry enough weight to form communities of their own and will simply join existing communities if they have already been agglomerated heavily, this would break our community detection capabilities. We have briefly touched on how agglomerative clustering could be provided in \cref{section:incr_hier_clustering}.

\subsection{Challenges}
\subsubsection{Rust}
Rust as a language proved to be a challenge when writing our implementation of the Differential Dataflow based Louvain method, this boiled down to the natural complexity that arose from dealing with the borrow checker that enforced affine types in a distributed/parallel setting with shared memory. In order to comply with the borrow checker, unfortunately, extra work was performed, however we used the properties of Differential Dataflow to our advantage here to keep these costs negligible after the initial computation was performed. 

\subsubsection{Differential Dataflow}
Unfortunately, the Rust implementation of Differential Dataflow is still quite undocumented and can be significantly difficult to understand due to the various mathematical abstractions it presents. This complexity contributed towards the difficulty regarding the evaluation of the performance characteristics of an algorithm of interest, while certain operations could be bounded by some algorithmic complexity in the worst case, at least theoretically, there was no realistically accessible method to evaluate the bounds of some such algorithm. In addition to these issues, there was missing documentation regarding the availability of some critical functionality such as the `explode' operator. This operator is crucial for high performance Differential Dataflow code as it typically can be used to replace the computationally expensive `reduce' operator. To understand how expensive the `reduce' operator is, we must first note that the reductions are performed in non-differential (ordinary) rust, resulting on two critical drawbacks, we are unable to incremental compute and our parallelism is limited artificially to the number of keys. While this statement is true, it is important to note that the `explode' operator is not as expressive as the `reduce' operator and for some logic such as reduction based on sorted values the `explode' operator simply is not applicable.

\subsection{Limitations}
Unfortunately as we have mentioned previously, there is no guarantee on convergence of our Louvain based method, in order to artificially induce convergence, we opted for a quick and easy method and that was to simply apply the identity function to a set of node, community pairs in the place where we would normally apply the clustering function. This is provides us with easy termination behaviour. 